\section*{Kurzfassung}
Helfen die nachrichtentechnischen Darstellungen (Dok.~288), Quarks und Neutrinos besser zu fassen?
Die ehrliche Antwort: \emph{als Sprache und Diagnose ja, als Herleitung nein} -- und für die beiden
Sektoren sehr unterschiedlich. Jedes Fermion-Triplett bekommt eine Magnitude-Koordinate (Koide $Q$,
bzw.\ $r=\sqrt{6Q-2}$) und eine Phasen-Koordinate. Gerechnet: nur die geladenen Leptonen sitzen am
speziellen Punkt $r=\sqrt2$ ($Q=2/3$); die Quarks liegen anderswo ($r\approx1{,}54$ bzw.\ $1{,}76$,
zudem skalenabhängig), und für Neutrinos lässt sich $Q$ gar nicht bilden. Der Gewinn ist eine
\emph{vereinheitlichende Karte}: Leptonen und Quarks sind betrags-dominiert (Hierarchie gross,
Mischung klein), Neutrinos phasen-dominiert (Betrag flach/entartet, Mischung gross -- PMNS/CKM-Gewicht
$0{,}55$ gegen $0{,}035$, Faktor~16). FFGFTs eigene Neutrino-Hypothese (gleiche Massen + Oszillation
aus geometrischer Phase) ist dann wörtlich der \emph{reine Allpass-Sektor}: Betrag flach, alles in der
Phase -- und sie wird falsifizierbar, weil die geometrische Phase die gemessenen $\Delta m^2$-Phasen
reproduzieren muss.

\section{Die Frage}
Dok.~288 hat gezeigt, dass der Z$_3$-Zirkulant-Massenoperator über die DFT konkrete Rechenhebel gibt.
Naheliegend ist: greifen dieselben Darstellungen auch bei den schwierigen Sektoren -- Quarks und
Neutrinos? Vorab nötig ist eine ehrliche Bestandsaufnahme, denn der Korpus behandelt diese Sektoren
nicht wie die sauberen geladenen Leptonen.

\section{Was der Korpus schon tut (pro Sektor)}
\begin{itemize}
  \item \textbf{Geladene Leptonen} (Dok.~006/046/258): der saubere Fall -- Foot-Koide-Form,
        $Q=2/3$, $r=\sqrt2$, Z$_3$-Zirkulant. Hier greift alles aus Dok.~288 unmittelbar.
  \item \textbf{Quarks} (Dok.~006/005): im Prinzip dieselbe $r_i/p_i$-Geometrie, aber die direkte
        geometrische Methode reicht nur $\sim1{,}2\%$; für Quarks braucht es die \emph{erweiterte
        fraktale Methode} mit QCD-Dynamik und \emph{ML-Kalibrierung an Lattice-Daten} (FLAG~2024).
        Also nicht rein geometrisch -- und Quarkmassen \emph{laufen} (skalen-/schemaabhängig).
  \item \textbf{Neutrinos} (Dok.~007/047): ein \emph{anderer} Mechanismus -- Photon-Analogie,
        $m_\nu=\tfrac{\xi_0^2}{2}m_e\approx4{,}54$~meV, alle drei Flavours (nahezu) gleich schwer,
        Oszillation aus \emph{geometrischen Phasen} statt Massendifferenzen. Im Korpus durchgehend
        als \textbf{spekulativ, ohne empirische Basis} markiert.
\end{itemize}

\section{Die (Q, r)-Koordinate: wo die Sektoren sitzen}
Über $Q=\tfrac13+\tfrac{r^2}{6}$ (also $r=\sqrt{6Q-2}$) bekommt jedes Triplett eine
Magnitude-Koordinate. Gerechnet (\texttt{koide\_q\_sektoren.py}, PDG-nahe Werte):

\begin{table}[htbp]
\centering
\caption{Magnitude-Koordinate der drei bildbaren Triplette. Quarkwerte sind skalenabhängig (MS-bar).}
\label{tab:qr}
\begin{tabular}{p{5.2cm}p{2.2cm}p{2.4cm}p{2.6cm}}
\toprule
\textbf{Sektor} & \textbf{Koide $Q$} & \textbf{$r=\sqrt{6Q-2}$} & \textbf{Lage} \\
\midrule
geladene Leptonen (e,\,$\mu$,\,$\tau$) & $0{,}6667$ & $\sqrt2=1{,}4142$ & speziell \\
up-Quarks (u,\,c,\,t)                  & $0{,}849$  & $1{,}759$         & nicht speziell \\
down-Quarks (d,\,s,\,b)                & $0{,}730$  & $1{,}543$         & nicht speziell \\
Neutrinos                              & ---        & ---               & kein Triplett \\
\bottomrule
\end{tabular}
\end{table}

Das ist schon die erste Antwort: \emph{nur die geladenen Leptonen sitzen am speziellen Punkt}
$r=\sqrt2$. Die Quark-$r$ sind nicht erkennbar besonders (und laufen mit der Skala); für Neutrinos
fehlt die absolute Skala und die Ordnung, sodass $Q$ nicht bildbar ist -- nur $\Delta m^2_{21}\approx
7{,}5\cdot10^{-5}$ und $\Delta m^2_{31}\approx2{,}5\cdot10^{-3}~\mathrm{eV}^2$ sind gemessen.

\section{Die Magnitude/Phase-Karte}
Die nachrichtentechnische Zerlegung (Betrag $=$ Massenspektrum, Phase $=$ Mischung/Oszillation) ordnet
die vier Sektoren in \emph{eine} Karte:
\begin{itemize}
  \item \textbf{Quarks}: starke Massenhierarchie (Betrag gross), kleine Mischung (CKM nahe Identität)
        -- \emph{betrags-dominiert}.
  \item \textbf{Geladene Leptonen}: Betrag an einem speziellen Wert ($r=\sqrt2$), kleine Phase
        ($\theta=2/9$) -- \emph{betrags-dominiert}.
  \item \textbf{Neutrinos}: flacher Betrag (entartet), grosse Mischung (PMNS weit von Identität)
        -- \emph{phasen-dominiert}.
\end{itemize}
Der CKM/PMNS-Kontrast lässt sich quantifizieren (\texttt{ckm\_pmns\_kontrast.py}, Off-Diagonal-Gewicht
der $|U|^2$-Matrix): CKM $0{,}035$ gegen PMNS $0{,}554$ -- ein Faktor $\sim16$. Die \enquote{kleine
gegen grosse Mischung} ist also genau ein \emph{Betrag-gegen-Phase}-Kontrast.

\section{Die Neutrinos als Allpass-Sektor}
Hier liegt der stärkste Punkt. FFGFTs eigene Hypothese -- gleiche Massen, Oszillation aus
geometrischer Phase -- ist in der Magnitude/Phase-Sprache wörtlich ein \emph{Allpass}: $|H|=1$
(flaches Betragsspektrum), nichttriviale Phase. Ein flaches Betragsspektrum kann keine Oszillation
kodieren; alles Physikalische sitzt in der Phase (\texttt{neutrino\_allpass.py}).

Das macht den Test \emph{konkret und falsifizierbar}. Die Standard-Oszillationsphase ist
$\Delta\phi_{ij}=1{,}267\,\Delta m^2_{ij}\,L/E$; bei T2K-artigem $L/E$ ($L=295$~km, $E=0{,}6$~GeV)
sind das $\Delta\phi_{31}\approx1{,}557$~rad und $\Delta\phi_{21}\approx0{,}047$~rad. FFGFTs
geometrische Phase (aus $\tilde T\cdot m=1$, Quantenzahlen $n,\ell,j$) \emph{muss} genau diese
$L/E$-abhängigen Phasen reproduzieren. Tut sie es nicht, ist die Entartungs-Hypothese widerlegt.

\section{Was das leistet -- und was nicht}
\begin{important}
\textbf{Leistung (diagnostisch).} Die Magnitude/Phase-Zerlegung erklärt die \emph{Arbeitsteilung}
zwischen den Sektoren: warum die Leptonen sauber sind (Betrag an speziellem Wert, kleine Phase), die
Quarks nicht (Betrag nicht speziell, extreme Hierarchie), die Neutrinos ganz anders (Betrag flach,
alles Phase). Das ist echtes Einordnungs-Wissen -- und es gibt der spekulativen Neutrino-Hypothese
eine scharfe, falsifizierbare Form.

\textbf{Grenze (P35).} Das ist bessere Sprache, eine vereinheitlichende Karte und ein schärferer Test
-- \emph{keine} neue Herleitung von Quark- oder Neutrinomassen. Die Quark-$r$ sind nicht speziell und
skalenabhängig, und der Quark-Sektor stützt sich im Korpus ohnehin auf QCD plus Lattice-Kalibrierung,
nicht auf reine Geometrie. Der Neutrino-Sektor bleibt spekulativ (eigene Markierung, Dok.~007) und in
Spannung zur sehr gut getesteten $\Delta m^2$-Deutung der Oszillation. Die saubere Allpass-Karte ist
eine bestechende Umformulierung, kein neuer Beleg.
\end{important}

\begin{conclusionBox}[Bilanz: eine Karte, kein neuer Sektor]
Die nachrichtentechnischen Darstellungen fassen Quarks und Neutrinos \emph{besser} -- aber im präzisen
Sinn: sie geben eine gemeinsame Magnitude/Phase-Koordinate, in der nur die geladenen Leptonen am
speziellen Punkt $r=\sqrt2$ sitzen, die Quarks betrags-dominiert daneben, und die Neutrinos als reiner
Phasen-/Allpass-Sektor. Der CKM/PMNS-Kontrast wird Betrag gegen Phase (Faktor~16). Und FFGFTs eigene
Neutrino-Hypothese erhält eine falsifizierbare Form: die geometrische Phase muss die gemessenen
$\Delta m^2$-Oszillationsphasen treffen. Das ist Diagnose und Test, keine Massenherleitung -- ehrlich
abgegrenzt.
\end{conclusionBox}
